\begin{textsample}{NEG Dim 4 – 2021\_c6 – Score 1.00 – 2021\_c6/enade\_2021\_hist\_1ca4f2173f.txt}  \label{ex:f4_neg_292}
CAPÍTULO 6 ANÁLISE TÉCNICA DA PROVA Este capítulo tem por objetivo apresentar o desempenho dos estudantes concluintes de História ( Licenciatura/Bacharelado ) no Enade/2021.
Para isso, foram calculadas as estatísticas básicas da prova em seu todo ( seção 6.1.1 ), bem como as estatísticas relacionadas aos Componentes de Formação Geral ( seção 6.1.2 ) e de Conhecimento Específico da Área ( seção 6.1.3 ).
Dadas as suas características, foram analisadas, em separado, as questões objetivas ( seção 6.2 ) e as questões discursivas ( seção 6.3 ).
Tomando -se como base as duas Questões Discursivas do Componente de Formação Geral, nas seções 6.3.1.6 a 6.3.1.8, são apresentados os resultados e os comentários sobre a correção das respostas em relação à Língua Portuguesa.
Para as questões objetivas, foram disponibilizados os Índices de Facilidade e de Discriminação Ponto-Bisserial, também em separado, para os Componentes de Formação Geral ( seção 6.2.1 e 6.2.2 ) e de Conhecimento Específico ( seção 6.2.3 e 6.2.4 ).
De cada componente, uma das questões foi escolhida para exemplificar a análise gráfica, relacionando -se as alternativas escolhidas pelos estudantes ( inclusive o gabarito ) com o número de acertos no componente.
Nos Anexos IA e IB, apresenta -se a íntegra da análise gráfica para todas as questões objetivas, respectivamente à Licenciatura e ao Bacharelado.
Para cada uma das questões discursivas, os conteúdos dos tipos mais comuns de respostas dos estudantes são apresentados e comparados com o padrão de resposta esperado ( ver Anexos IXA e IXB com o padrão de respostas ).
Nas tabelas, constam as seguintes estatísticas das notas19 : média do desempenho na prova, erro padrão da média, desvio padrão, nota mínima, mediana e nota máxima para cada um de seus componentes.
Tais estatísticas contemplam o total de estudantes concluintes da Área de História ( Licenciatura/Bacharelado ) inscritos e presentes à prova do Enade/2021, tendo em vista agregações, ou por Grandes Regiões e o país como um todo, segundo Habilitação e Modalidade de Ensino, ou por Categoria Administrativa, Organização Acadêmica da IES.
Em relação aos gráficos de distribuição de notas, o intervalo considerado foi de 10 unidades, aberto à esquerda e fechado à direita, com exceção do primeiro intervalo, [ 0 ; 10 ], fechado em ambos os extremos.
Para os gráficos de distribuição das notas das questões discursivas, foram consideradas mais duas categorias : questão em branco20 e nota zero. 19 Uma definição dessas estatísticas pode ser encontrada no Glossário. 20 Nesse grupo estão incluídas também as respostas classificadas como nulas ou desconsideradas. 6.1 ESTATÍSTICAS BÁSICAS DA PROVA Esta seção apresenta estatísticas selecionadas e histogramas da nota geral ( 6.1.1 ) e de cada componente : Formação Geral ( 6.1.2 ) e Conhecimento Específico ( 6.1.3 ).
São também apresentadas estatísticas selecionadas de subpopulações, caracterizadas por Grande Região, Categoria Administrativa, Organização Acadêmica e Modalidade de Ensino. 6.1.1 ESTATÍSTICAS BÁSICAS GERAIS Na Tabela 6.1, são apresentadas as Estatísticas Básicas da prova, por Grande Região segundo Habilitação e Modalidade de Ensino, dos estudantes concluintes de História ( Licenciatura/Bacharelado ).
As análises serão feitas primeiramente para os cursos de Licenciatura e depois para os de Bacharelado.
Cada habilitação é desagregada internamente em cursos Presenciais e cursos a Distância.
A Média das notas da prova, como um todo, dos alunos de Licenciatura ( nas seções seguintes serão analisados os Componentes de Formação Geral e de Conhecimento Específico ), foi 42,6, sendo que os estudantes da região Sul obtiveram a Média mais baixa ( 39,9 ), e os da região Sudeste obtiveram a Média mais alta ( 45,0 ).
As demais Médias foram : 40,7 na região Norte, 43,0, na região Nordeste, e 42,3, na região Centro-Oeste.
O Desvio padrão para o Brasil, como um todo, foi 16,4, sendo o maior Desvio padrão encontrado na região Sudeste ( 17,1 ), e o menor, na região Norte ( 14,8 ), indicando uma dispersão um pouco menor das notas dessa última região.
A região que obteve a maior nota Máxima foi a Sul ( 88,5 ), ao passo que a região que atingiu a menor nota Máxima foi a Norte ( 75,0 ).
A Mediana do Brasil, como um todo, foi 41,3, sendo a maior Mediana obtida na região Sudeste ( 44,4 ), e a menor obtida na região Sul ( 37,5 ).
A nota Mínima do Brasil foi zero, obtida em quatro das cinco regiões.
A exceção foi a nota Mínima da região Norte, foi 4,7.
Considerando -se as notas, como um todo dos alunos de Licenciatura segundo Grande Região, observa -se que não existe diferença estatisticamente significativa ao nível de 95 \% entre as Médias obtidas nas regiões Norte e Sul, as duas menores Médias e entre as Médias das regiões Nordeste e Centro-Oeste.
A Média das notas para os alunos de cursos Presenciais de Licenciatura, foi 47,4, sendo que os alunos da região Norte obtiveram a Média mais baixa ( 40,7 ), e os da região Sul obtiveram a Média mais alta ( 52,4 ).
As demais Médias foram : 44,4 na região Nordeste, 51,7 na região Sudeste e 43,0 na região Centro-Oeste.
O Desvio padrão das Licenciaturas Presenciais foi 16,2, sendo o menor obtido na região Norte 14,8 e o maior na região Centro-Oeste ( 16,5 ).
Oeste ( 16,5 ).
A maior Mediana foi a da região Sudeste ( 53,5 ) e a menor foi a da região Norte ( 40,6 ).
A maior nota Máxima foi 88,5, obtida na região Sul.
As notas Mínimas foram de zero em quatro regiões ( Nordeste, Sudeste, Sul e Centro-Oeste ) e na região Norte foi 4,7.
Na Modalidade de Educação Presencial de Licenciatura as diferenças entre as duas maiores Médias, obtidas nas regiões Sul e Sudeste, não são estatisticamente significativas.
Também não são estatisticamente significativas as diferenças entre as Médias das regiões Nordeste e Centro-Oeste.
As diferenças entre quaisquer outros pares de Médias são significativas.
Na Modalidade a Distância não houve participação de cursos oferecidos na região Norte.
A Média das notas da prova dos alunos de cursos a Distância de Licenciatura, foi 36,9, mais baixa do que os alunos de cursos Presenciais ( 47,4 ).
A Média mais alta foi obtida na região Sudeste ( 38,2 ) e a mais baixa na região Nordeste ( 34,8 ).
O Desvio padrão das notas dos cursos a Distância foi 14,8, sendo o maior obtido na região Sudeste ( 15,5 ) e o menor na região Nordeste ( 13,3 ).
Para essa combinação de Modalidade e Habilitação, a nota Máxima foi 88,1, obtida, por pelo menos, um estudante da região Sul.
As demais notas Máximas foram : 80,0 na região Nordeste, 86,3 na região Sudeste e 83,0 na região Centro-Oeste.
A Mediana dos cursos a Distância do Brasil como um todo foi 34,7.
A nota Mínima foi 3,8 na região Nordeste e 8,2 na região Centro-Oeste.
Nas regiões Sudeste e Sul a nota Mínima foi zero.
Cumpre notar que ainda que as notas dos cursos Presenciais sejam comparáveis com as dos cursos a Distância dentro de uma mesma Habilitação, no entanto, não é possível comparar as notas obtidas pelos alunos de Licenciatura e de Bacharelado, já que as provas são, em parte, diferentes.
A Média das notas da prova dos alunos de Bacharelado ( nas seções seguintes serão analisados os Componentes de Formação Geral e de Conhecimento Específico ), foi 53,9, sendo que os estudantes da região Norte obtiveram a Média mais baixa ( 35,6 ), e os da região Centro-Oeste obtiveram a Média mais alta ( 57,1 ).
As demais Médias foram : 51,1, na região Nordeste, 54,9, na região Sudeste, e 53,2, na região Sul.
O Desvio padrão para o Brasil, como um todo, foi 15,3, sendo o maior Desvio padrão encontrado na região Centro-Oeste ( 16,2 ), e o menor, na região Norte ( 12,9 ), indicando uma dispersão um pouco menor das notas dessa última região.
A região que obteve a maior nota Máxima foi a Sudeste ( 88,6 ), ao passo que a região que atingiu a menor nota Máxima foi a Norte ( 57,8 ).
A Mediana do Brasil, como um todo, foi 55,8, sendo a maior Mediana obtida na região Sudeste ( 56,4 ), e a menor obtida na região Norte ( 36,5 ).
A nota Mínima do Brasil foi obtida na região Sul, onde foi zero.
A maior nota Mínima, entre as cinco Grandes Regiões, foi 12,4, na região Sudeste, as demais regiões as notas Mínimas foram : 7,4 na região Norte, 9,3 na região Nordeste, e 7,5 na região Centro-Oeste.
Considerando -se as notas de todo o Brasil dos alunos de Bacharelado, segundo Grande Região, observa -se que existe diferença estatisticamente significativa ao nível de 95 \% apenas entre a Média obtida pela região Norte e as Médias das demais regiões.
A Média das notas da prova como um todo para os alunos de cursos Presenciais de Bacharelado, foi 55,2, sendo que os alunos da região Norte obtiveram a Média mais baixa ( 35,6 ), e os da região Centro-Oeste obtiveram a Média mais alta ( 57,1 ).
As demais Médias foram : 51,1 na região Nordeste, 56,4 na região Sudeste e 54,7 na região Sul.
A maior nota Máxima dos cursos Presenciais de Bacharelado foi a da região Sudeste ( 88,6 ) e a menor foi a da região Norte ( 57,8 ).
A maior Mediana foi obtida na região Sul ( 57,5 ), e a menor obtida na Norte ( 36,5 ).
As notas Mínimas nos cursos Presenciais foram : 7,4, na região Norte, 9,3, na região Nordeste, 12,4 na região Sudeste, zero na região Sul, e, 7,5 na região Centro-Oeste.
Considerando -se as notas segundo Grande Região, é estatisticamente significativa ao nível de 95 \% a diferença entre a Média da região Norte e todas as demais.
O mesmo ocorre com a diferença entre a Média da região Nordeste e a da região Sudeste.
Somente nas regiões Sudeste e Sul foram oferecidos cursos a Distância de Bacharelado.
A Média foi 45,8 e o Desvio padrão foi 16,2.
Estas estatísticas foram maiores na região Sul do que na região Sudeste.
Para essa combinação de Modalidade e Habilitação, a nota Máxima foi 84,4, obtida na região Sul.
As Medianas foram 39,8 na região Sudeste e 51,6 na região Sul.
A nota Mínima foi 14,9 na região Sudeste e 11,0 na região Sul.
O comportamento das notas dos estudantes de todo o Brasil pode ser observado no Gráfico 6.1, que apresenta um histograma com a distribuição das mesmas, segundo a Habilitação.
Essa é uma distribuição unimodal para os alunos de cursos de Licenciatura e de Bacharelado.
A moda da Licenciatura ocorre no intervalo ( 30 ; 40 ] e do Bacharelado no intervalo ( 60 ; 70 ], com respectivamente 21,6 \% e 24,2 \% das respostas da Habilitação.
Na Tabela 6.2, são apresentadas informações referentes à Média da nota final, desagregadas por Categoria Administrativa, por Organização Acadêmica, segundo a Habilitação.
A Média das notas da prova, como um todo, dos alunos de Licenciatura, foi 42,6 ( Tabela 6.1 ).
Em relação a Categoria Administrativa, os estudantes dos cursos de Licenciatura das IES Públicas obtiveram Média mais alta ( 46,8 ) que a Média nacional, e os das IES Privadas obtiveram Média mais baixa ( 38,7 ).
Observa -se que existe diferença estatisticamente significativa entre as Médias das notas das IES Públicas e as das Privadas.
A diferença entre as Médias das regiões Sudeste e Sul ( 5,1 ), a maior e a menor Média, é inferior à diferença entre a das IES Públicas e Privadas ( 8,1 ), caracterizando que a diversidade administrativa é ligeiramente maior do que a diversidade regional.
O Desvio padrão para as IES Públicas ( 16,4 ) igual ao do Brasil, como um todo, indicando que a dispersão nesta Categoria Administrativa é igual a do total de participantes de Licenciatura.
No tocante à Organização Acadêmica, os CEFET/IF ( 54,5 ) e as Universidades ( 44,1 ) obtiveram Média mais alta que a nacional ( 42,6 ).
A Média dos Centros Universitários ( 38,8 ) e a das Faculdades ( 34,0 ) foram menores do que a nacional.
Constata -se que existe diferença estatisticamente significativa, ao nível de 95 \%, entre as Médias das Organizações Acadêmicas.
A Média das notas da prova, como um todo, dos alunos de Bacharelado, foi 53,9 ( Tabela 6.1 ).
Em relação a Categoria Administrativa, os estudantes de Bacharelado das IES Públicas obtiveram Média mais alta ( 55,3 ) que a Média nacional, e os das IES Privadas obtiveram Média mais baixa ( 48,2 ).
Observa -se que existe diferença estatisticamente significativa entre as Médias das notas das IES Públicas e as das Privadas.
A diferença entre as Médias das regiões Centro-Oeste e Norte ( 21,5 ), a maior e a menor Média, é bem superior à diferença entre a das IES Públicas e Privadas ( 7,1 ), caracterizando -se uma maior diversidade regional do que administrativa.
O Desvio padrão para as IES Públicas ( 14,9 ) foi inferior ao do Brasil, como um todo ( 15,3 ), e as IES Privadas ( 15,9 ), ligeiramente superior, indicando uma dispersão um pouco maior das notas nessa Categoria Administrativa.
No tocante a Organização Acadêmica, não houve participação de estudantes de Bacharelado em História de cursos oferecidos por Faculdades e CEFET/IF.
As Universidades obtiveram Média mais alta que a nacional ( 54,3 ).
A Média dos Centros Universitários foi menor do que a nacional ( 49,6 ).
Constata -se que existe diferença estatisticamente significativa, ao nível de 95 \%, entre as Médias dos dois tipos de Organização Acadêmica representados. 6.1.2 ESTATÍSTICAS BÁSICAS NO COMPONENTE DE FORMAÇÃO GERAL Na Tabela 6.3, são apresentadas as Estatísticas Básicas em relação ao componente da prova que avalia a Formação Geral dos estudantes concluintes.
A nota de Formação Geral é obtida levando -se em conta os acertos nas oito questões objetivas e a nota Média obtida nas Questões Discursivas 1 e 2, para ambas Habilitações, sendo a nota de cada questão discursiva obtida levando -se em conta 80 \% da nota relativa ao conteúdo e 20 \% da nota relativa ao desempenho linguístico.
Os estudantes de Licenciatura, de todo o Brasil, obtiveram desempenho médio de 36,9.
Quanto à variabilidade, o Desvio padrão das notas dos estudantes de Licenciatura do Brasil foi 15,0.
A maior Média foi obtida na região Sudeste ( 38,4 ) e a menor, nas regiões Norte e Sul ( 35,8 ).
As demais Médias foram : 36,4, na região Nordeste e 36,8, na região Centro-Oeste.
Já o maior Desvio padrão foi obtido na região Centro-Oeste ( 15,4 ) e o menor, na região Norte ( 13,7 ).
Os demais desvios foram : 14,2, na região Nordeste e 15,3, nas regiões Sudeste e Sul.
A maior nota no Componente de Formação Geral da Licenciatura em História na prova do Enade/2021 foi 89,8, obtida por, pelo menos, um estudante na região Sul.
A menor nota Máxima foi obtida na região Norte ( 70,4 ).
A Mediana do Brasil, como um todo, foi 36,7, sendo a menor Mediana encontrada na região Sul ( 34,9 ), e a maior, na região Sudeste ( 38,1 ).
A nota Mínima nessa parte da prova foi zero em todas as regiões.
Considerando -se as notas, segundo Grande Região, observa -se que não existe diferença estatisticamente significativa entre as Médias das notas do Componente de Formação Geral obtidas em quatro regiões ( Norte, Nordeste, Sul e Centro-Oeste ).
As diferenças entre as Médias da região Sudeste em relação às demais regiões são estatisticamente significativas.
Os alunos de cursos Presenciais de Licenciatura obtiveram desempenho médio de 39,6.
Quanto à variabilidade, o Desvio padrão das notas do conjunto dos estudantes das Licenciaturas Presenciais foi 14,7.
A maior Média foi obtida na região Sul ( 43,0 ), e a menor, na região Norte ( 35,8 ).
As demais Médias foram : 37,2 na região Nordeste, 42,7 na região Sudeste e 37,4 na região Centro-Oeste.
O maior Desvio padrão foi obtido na região Centro-Oeste ( 15,4 ), e o menor, na região Norte ( 13,7 ).
Os demais desvios foram : 14,0 na região Nordeste, 14,6 na região Sudeste e 14,8 na região Sul.
A nota Máxima no Componente de Formação Geral da prova de alunos de cursos Presenciais de Licenciatura foi 89,8, obtida por, pelo menos, um aluno na região Sul.
A menor nota Máxima foi obtida na região Norte ( 70,4 ).
A Mediana do conjunto de Licenciaturas Presenciais foi 39,4, sendo a menor Mediana encontrada na região Norte ( 35,4 ), e a maior, encontrada na região Sul ( 43,2 ).
A nota Mínima nesta parte foi zero em todas as regiões.
Considerando -se as notas segundo Grande Região das Licenciaturas Presenciais, no Componente de Formação Geral, observa -se que não existe diferença estatisticamente significativa entre as duas maiores Médias, obtidas nas regiões Sul e Sudeste.
Da mesma forma, não são significativas as diferenças entre as três Médias menores das notas, obtidas nas regiões Norte, Nordeste e Centro-Oeste.
Na região Norte não houve participação de estudantes de cursos a Distância de Licenciatura.
A Média das notas do Componente de Formação Geral dos alunos de Licenciatura a Distância foi 33,7, um pouco mais baixa do que a dos cursos Presenciais ( 39,6 ), o mesmo ocorrendo com a Mediana ( 32,9 contra 39,4 para os cursos Presenciais ).
A nota Mínima também foi zero para todas as regiões.
Para essa combinação de Modalidade e Habilitação, o Desvio padrão foi 14,8, e a nota Máxima foi 84,2.
Os estudantes de Bacharelado, de todo o Brasil, obtiveram desempenho médio de 44,1 ( Tabela 6.3 ).
Quanto à variabilidade, o Desvio padrão das notas dos estudantes de Bacharelado do Brasil, como um todo, foi 15,1.
A maior Média foi obtida nas regiões Sudeste e Centro-Oeste ( 45,4 ), e a menor, na região Norte ( 28,0 ).
As demais Médias foram : 39,2, na região Nordeste ; 43,2 e na região Sul.
Já o maior Desvio padrão foi obtido na região Centro-Oeste ( 18,4 ) e o menor, na região Nordeste ( 13,9 ).
Na região Norte o Desvio padrão foi 15,0, na região Sudeste foi 14,4 e na região Sul foi 15,8.
A maior nota no Componente de Formação Geral da prova do Enade/2021 foi 84,4, obtida por, pelo menos, um estudante na região Sudeste.
A menor nota Máxima foi obtida na região Norte ( 53,3 ).
A Mediana do Brasil, como um todo, foi 45,0, sendo a menor Mediana encontrada na região Norte ( 24,6 ), e a maior, na região Centro-Oeste ( 47,9 ).
A nota Mínima nessa parte foi diferente de zero em três das cinco Grandes Regiões e nas demais foi : Nordeste com 14,9 e Centro-Oeste com 7,5.
Considerando -se as notas, segundo Grande Região, observa -se que existe diferença estatisticamente significativa entre a Média das notas do Componente de Formação Geral obtida na região Norte e as Médias das demais regiões.
Como não houve participação de cursos de Bacharelado em História oferecidos na Modalidade a Distância nas regiões Norte, Nordeste e Centro-Oeste, as estatísticas das notas dos alunos de cursos Presenciais de Bacharelado nessas regiões são as mesmas apresentadas para o total de cursos de Bacharelados participantes do Enade em 2021.
A desagregação por Modalidade de Ensino ocorre apenas nas regiões Sudeste e Sul.
A maior Média foi obtida na região Sul ( 43,6 ) e a menor na região Sudeste ( 36,0 ).
A região que obteve a maior nota Máxima foi a Sudeste ( 83,8 ) e a nota Máxima da região Sul foi 80,6.
A Mediana do conjunto dos Bacharelados a Distância foi 40,8, sendo a Mediana da região Sudeste ( 37,5 ) menor do que a da região Sul ( 46,0 ).
As notas Mínimas foram 7,5 e 3,0, respectivamente nas regiões Sudeste e Sul.
Considerando -se as notas dessas duas regiões, observa -se que existe diferença estatisticamente significativa ao nível de 95 \% entre as suas Médias.
Os estudantes de Bacharelado na Modalidade Presencial, de todo o Brasil, obtiveram desempenho médio de 44,7 ( Tabela 6.3 ).
Quanto à variabilidade, o Desvio padrão das notas dos estudantes de cursos Presenciais de Bacharelado do Brasil, como um todo, foi 15,0.
A maior Média foi obtida na região Sudeste ( 46,3 ), e a menor, na região Norte ( 28,0 ).
As demais Médias foram : 39,2, na região Nordeste, 43,1 e na região Sul e 45,4 na região Centro-Oeste.
Já o maior Desvio padrão foi obtido na região Centro-Oeste ( 18,4 ) e o menor, na região Nordeste ( 13,9 ).
Na região Norte o Desvio padrão foi 15,0, na região Sudeste foi 14,0 e na região Sul foi 16,2.
A maior nota Máxima foi atingida na região Sudeste ( 84,4 ) e a menor na região Norte ( 53,3 ), como nos resultados para o Bacharelado Presencial como um todo.
Da mesma forma, os resultados de notas Mínimas foram os mesmos de antes da desagregação : a maior nota Mínima foi a da região Nordeste ( 14,9 ), a segunda foi a da região Centro-oeste ( 7,5 ), e a das demais regiões foi zero.
O Gráfico 6.2 propicia a avaliação do desempenho dos estudantes do Componente de Formação Geral, a partir do histograma da distribuição das notas correspondentes.
A distribuição é unimodal para as notas de Licenciatura com moda em ( 30 ; 40 ], mesmo intervalo modal da distribuição de notas da prova, como um todo ( Gráfico 6.1 ).
Para as notas do Bacharelado, a moda ocorre no intervalo ( 40 ; 50 ], bem inferior ao modal da distribuição de notas da prova, como um todo ( Gráfico 6.1 ).
Nota -se, ainda, que no Gráfico 6.2, as notas apresentam dispersão menor, quando comparadas a do Gráfico 6.1 ( distribuição das notas da prova ), confirmado pela comparação dos desvios padrões.
Em Formação Geral, para a Licenciatura, o desvio diminui de 16,4 para a nota da prova, como um todo ( Tabela 6.1 ), para 15,0 para o Componente de Formação Geral ( Tabela 6.3 ).
Para os Bacharelados, diminui de 15,3 para a nota da prova, como um todo ( Tabela 6.1 ), para 15,1 para o Componente de Formação Geral ( Tabela 6.3 ).
Na Tabela 6.4, são apresentadas as informações referentes ao desempenho dos concluintes do Componente de Formação Geral, em diferentes agregações : Categoria Administrativa, Organização Acadêmica e Habilitação.
Primeiramente, com respeito aos alunos de Licenciatura, observa -se que existe diferença estatisticamente significativa entre as Médias dos tipos de Categoria Administrativa : a maior Média obtida por estudantes de Licenciatura de IES Públicas ( 39,1 ) e a menor, pelos de IES Privadas ( 34,9 ).
Nota -se que, também, há diferença estatisticamente significativa entre as Médias dos tipos de Organização Acadêmica : 37,7 para as Universidades, 34,9 para os Centros Universitários, 32,2 para as Faculdades, a menor Média, e 43,1 para os CEFET/IF, a maior.
Dentre os alunos de Bacharelado, observa -se que existe diferença estatisticamente significativa entre as Médias dos tipos de Categoria Administrativa : a maior Média obtida por estudantes de Bacharelado de IES Públicas ( 44,8 ) e a menor, pelos de IES Privadas ( 41,3 ).
Nota -se que, não há diferença estatisticamente significativa entre as Médias dos dois tipos de Organização Acadêmica que tiveram estudantes de Bacharelado em História participando do Enade/2021, o maior valor para os Universidades ( 44,2 ) e o menor, para os Centros Universitários ( 43,1 ). 6.1.3 ESTATÍSTICAS BÁSICAS DO COMPONENTE DE CONHECIMENTO ESPECÍFICO Na Tabela 6.5, são apresentadas as Estatísticas Básicas referentes ao Componente de Conhecimento Específico da Área de História ( Licenciatura/Bacharelado ) por Grande Região, segundo Habilitação e Modalidade de Ensino.
A nota do componente de Conhecimentos Específicos leva em conta as notas da parte objetiva e a Média das notas obtidas nas três questões discursivas, parte discursiva da Área.
A Média do desempenho dos estudantes de Licenciatura do Brasil, como um todo, é 44,5.
A maior Média é obtida na região Sudeste ( 47,1 ), e a menor, na região Sul ( 41,2 ).
As demais Médias são : 42,4 na região Norte, 45,2 na região Nordeste e 44,1 na região Centro-Oeste.
Quanto à variabilidade das notas, o Desvio padrão do Brasil, como um todo para os alunos de Licenciatura, é 18,7, sendo o maior Desvio padrão observado na região Sudeste ( 19,4 ) e o menor, na região Norte ( 17,0 ).
Os demais desvios são : 17,9 na região Nordeste, 18,3 na região Sul e 18,4 na região Centro-Oeste.
A Mediana das notas dos estudantes de Licenciatura de todo o Brasil é 42,5.
A maior Mediana ocorre na região Sudeste ( 46,3 ), e a menor, na região Sul ( 38,5 ).
As demais Medianas são : 42,4 na região Norte, 43,7 na região Nordeste e 42,4 na região Centro-Oeste.
A nota Máxima do Brasil, como um todo, é 98,8, sendo obtida por pelo menos um estudante de Licenciatura na região Sudeste.
As demais notas Máximas são : 86,0 na região Norte, 96,3 na região Nordeste, 97,5 na região Sul e 92,4 na região Centro-Oeste.
A nota Mínima é zero obtida por pelo menos um estudante em todas as regiões.
Observa -se que existe diferença estatisticamente significativa entre as regiões Norte, Nordeste e Sudeste na análise das Médias das notas do Componente de Conhecimento Específico, bem como entre a região Sul e as regiões Nordeste, Sudeste e Centro-Oeste, entre a região Sudeste e Centro-Oeste.
Já para os alunos dos cursos Presenciais de Licenciatura, a Média do desempenho de Conhecimento Específico é 49,9.
A maior Média ocorre na região Sul ( 55,6 ), e a menor, na região Norte ( 42,4 ).
As demais Médias são : 46,8 na região Nordeste, 54,7 na região Sudeste e 44,9 na região Centro-Oeste.
Quanto à variabilidade das notas, o Desvio padrão do conjunto das Licenciaturas Presenciais é 18,6, sendo o maior Desvio padrão observado na região Centro-Oeste ( 18,8 ), e o menor, na região Norte ( 17,0 ).
Os demais desvios são : 17,8 na região Nordeste, 18,3 na região Sudeste e 18,0 na região Sul.
A Mediana das notas dos estudantes de cursos Presenciais de Licenciatura de todo o Brasil é 50,1.
A maior Mediana ocorre na região Sul ( 57,8 ), e a menor, na região Norte ( 42,4 ).
As demais Medianas são : 46,3 na região Nordeste, 56,6 na região Sudeste e 43,7 na região Centro-Oeste.
A nota Máxima das Licenciaturas Presenciais em História é 98,8, sendo obtida por, pelo menos, um aluno na região Sudeste.
As demais notas Máximas são : 86,0 na região Norte, 96,3 na região Nordeste, 97,5 na região Sul e 92,4 na região Centro-Oeste.
A nota Mínima é zero obtida por pelo menos um estudante em todas as regiões.
Considerando -se as Licenciaturas oferecidas presencialmente, observa -se que não existe diferença estatisticamente significativa entre as Médias das notas do Componente de Conhecimento Específico das regiões Norte e Centro-Oeste, e entre as Médias das regiões Sudeste e Sul.
Cursos a Distância de Licenciatura em História não foram ofertados na região Norte.
A Média do desempenho de Conhecimento Específico dos alunos dos cursos a Distância de Licenciatura é 38,0, menor do que a dos cursos Presenciais.
A maior Média ocorre na região Sudeste ( 39,5 ) e a menor na região Nordeste ( 35,7 ).
Quanto à variabilidade das notas, o Desvio padrão é 16,7, sendo o maior Desvio padrão observado na região Sudeste ( 17,5 ).
A nota Máxima, 95,0, foi obtida por, pelo menos um aluno na região Sul.
Apenas na região Centro-Oeste a nota Mínima é diferente de zero ( 6,3 ).
Tanto a região Sudeste quanto a região Centro-Oeste têm as maiores Medianas, 37,3.
A Média do desempenho dos estudantes de Bacharelado do Brasil, como um todo, é 57,1.
A maior Média ocorre na região Centro-Oeste ( 61,0 ), e a menor, na região Norte ( 38,2 ).
As demais Médias são : 55,0 na região Nordeste, 58,1 na região Sudeste e 56,5 na região Sul.
Quanto à variabilidade das notas, o Desvio padrão do Brasil, como um todo para os alunos de Bacharelado, é 17,4, sendo o maior Desvio padrão observado na região Sul ( 18,1 ) e o menor, na região Norte ( 14,3 ).
Os demais desvios são : 17,8 na região Nordeste, 16,7 na região Sudeste e 17,1 na região Centro-Oeste.
A Mediana das notas dos estudantes de Bacharelado de todo o Brasil é 59,3.
A maior Mediana ocorre na região Centro-Oeste ( 62,3 ), e a menor, na região Norte ( 39,5 ).
As demais Medianas são : 59,2 na região Nordeste, 60,4 na região Sudeste e 59,2 na região Sul.
A nota Máxima do Brasil, como um todo, é 95,1, sendo obtida por, pelo menos, um estudante de Bacharelado na região Sudeste.
As demais notas Máximas são : 69,1 na região Norte, 84,0 na região Nordeste, 91,4 na região Sul e 90,1 na região Centro-Oeste.
A nota Mínima é zero na região Sul.
A maior nota Mínima ocorre na região Sudeste ( 9,9 ).
Observa -se que existe diferença estatisticamente significativa entre as Médias, das notas do Componente de Conhecimento Específico, da região Norte e as demais regiões.
Já para os alunos de cursos Presenciais de Bacharelado, a Média do desempenho de Conhecimento é ligeiramente maior, 58,7.
A maior e a menor Média são respectivamente nas regiões Centro-Oeste ( 61,0 ) e Norte ( 38,2 ).
Nas regiões Sudeste e Sul, que têm a participação de cursos a Distância, a Média desagregada dos cursos Presenciais aumenta de 58,1 para 59,7 e de 56,5 para 58,6, respectivamente.
O Desvio padrão dos cursos Presenciais da região Sudeste, passa de 16,7 para 15,9, e com a Mediana, de 60,4 para 60,5.
E o Desvio padrão dos cursos Presenciais da região Sul, passa de 18,1 ( mais alto ) para 17,8, e com a Mediana, de 59,2 para 60,5.
As notas Máximas e as notas Mínimas obtidas em todas as regiões para os Bacharelados Presenciais, na maioria dos casos, foram as mesmas apresentadas para os Bacharelados considerando -se o agregado das duas Modalidades, exceto a Máxima na região Sul que nesse caso é 88,9.
Existe diferença estatisticamente significativa entre as Médias das notas do Componente de Conhecimento Específico da região Norte com as demais regiões.
Cursos a Distância de Bacharelado em História são ofertados nas regiões Sudeste e Sul.
As estatísticas calculadas por região obtêm valores mais elevados na região Sul do que na região Sudeste.
As Médias são 52,1 e 41,7 e os Desvios padrões são 18,2 e 16,4.
As notas Máxima e Mínima são, respectivamente, ( 91,4 e 77,7 ) e ( 12,3 e 11,1 ).
E as Medianas são 53,0 e 40,6.
Assim como nos Gráficos 6.1 e 6.2, no Gráfico 6.3, apresentado a seguir, é possível fazer uma avaliação do desempenho de concluintes de História, segundo a Habilitação, em relação ao Componente de Conhecimento Específico, com um histograma da distribuição das notas correspondentes.
Essa também é uma distribuição unimodal para ambas as Habilitações.
Para os estudantes de Licenciatura, o grupo modal é o intervalo de notas ( 30 ; 40 ] tanto para o Componente de Conhecimento Específico ( Gráfico 6.3 ) quanto para a Formação Geral ( Gráfico 6.2 ) e a nota geral da prova ( Gráfico 6.1 ).
Para os estudantes de Bacharelado, o grupo modal para o Componente de Conhecimento Específico ( Gráfico 6.3 ) é o mesmo da prova como um todo ( Gráfico 6.1 ), o intervalo de notas ( 60 ; 70 ], sendo esse superior ao da Formação Geral ( Gráfico 6.2 ), que é o intervalo ( 40 ; 50 ].
Na Tabela 6.6, apresenta -se uma comparação dos resultados em relação a Categoria Administrativa, a Organização Acadêmica e segundo a Habilitação, agora levando -se em conta o desempenho de estudantes do Componente de Conhecimento Específico da prova.
No que se refere a Organização Acadêmica, para os alunos de Licenciatura, a maior Média é a dos CEFET/IF ( 58,2 ), vindo a seguir a das Universidades ( 46,2 ) ; o valor obtido para os Centros Universitários ( 40,1 ) e para as Faculdades ( 34,5 ) é inferior ao Brasil ( 44,5 ).
O maior Desvio padrão, e o único superior ao do Brasil ( 18,7 ), como um todo, é o das Universidades ( 18,9 ).
As Universidades obtêm a maior nota Máxima ( 98,8 ).
Os Centros Universitários obtêm nota Máxima 95,0, e as Faculdades e os CEFET/IF, 85,9.
As Medianas são : 45,0 nas Universidades, 37,3 nos Centros Universitários, 30,9 nas Faculdades e 59,1 nos CEFET/IF.
A nota Mínima é zero para metade das Organizações Acadêmicas, Faculdades ( 1,2 ) e CEFET/IF ( 21,7 ).
Observa -se que as diferenças, entre a Média das notas das IES Públicas e Privadas, e também entre a Média das notas de todos os tipos de Organização Acadêmica, são estatisticamente significativas, ao nível de 95 \%, no Componente de Conhecimento Específico.
Em relação a Categoria Administrativa, verifica -se que as IES Públicas apresentam sempre valores mais elevados que as IES Privadas.
A Média das IES Públicas é 49,4, o Desvio padrão é 18,8, a nota Máxima é 98,8, a Mínima é zero e a Mediana é 50,0.
No que se refere a Organização Acadêmica, para os alunos de Bacharelado, apenas as Universidades e os Centros Universitários têm cursos avaliados.
A Média, a Mediana e a nota Máxima são mais elevadas para as Universidades do que para os Centros Universitários e o Desvio padrão e a Mínima têm comportamento contrário.
Em relação a Categoria Administrativa o mesmo ocorre, a Média, a Mediana e a Máxima são mais elevadas para as IES Públicas do que para as IES Privadas e o Desvio padrão e a Mínima têm comportamento contrário.
Observa -se que as diferenças, entre a Média das notas de IES Públicas e Privadas, e entre a Média das notas entre as Universidades e os Centros Universitários, são estatisticamente significativas, ao nível de 95 \%, no Componente de Conhecimento Específico. 6.2 ANÁLISE DAS QUESTÕES OBJETIVAS Esta seção apresenta estatísticas selecionadas e histogramas das Questões Objetivas de Formação Geral para Licenciatura ( 6.2.1 ) e para Bacharelado ( 6.2.2 ) e de Conhecimento Específico para Licenciatura ( 6.2.3 ) e para Bacharelado ( 6.2.4 ).
São, também, apresentadas e comparadas as Médias das subpopulações caracterizadas por Grande Região. 21 21 As fórmulas utilizadas para o coeficiente de correlação ponto-bisserial e para o índice de facilidade estão apresentadas no capítulo 1. 6.2.1 COMPONENTE DE FORMAÇÃO GERAL – HISTÓRIA ( LICENCIATURA ) Na Tabela 6.7a, são apresentadas as Estatísticas Básicas relativas às oito Questões Objetivas do Componente da prova que abrange a Formação Geral dos estudantes de Licenciatura.
A Média do Brasil é 42,6.
A menor Média é encontrada na região Sul ( 40,5 ), e a maior, na região Sudeste ( 44,8 ).
As demais Médias são : 41,3 na região Norte, 42,4 na região Nordeste e 42,5 na região Centro-Oeste.
O Desvio padrão do Brasil é 19,9, sendo o maior Desvio padrão encontrado na região Centro-Oeste ( 20,2 ), e o menor Desvio padrão encontrado na região Norte ( 18,4 ).
Os demais Desvios padrões são : 19,3 na região Nordeste, 20,1 na região Sudeste e 20,0 na região Sul.
As Medianas do Brasil, como um todo, e de quatro regiões são iguais a 37,5 ; exceto na região Sudeste que o valor é 50,0.
A nota Máxima 100,0 é alcançada em quatro das cinco Grandes Regiões ; na região Norte a Máxima é 87,5.
As notas Mínimas ( 0,0 ) são iguais para todas as regiões.
Na Tabela 6.8a são apresentados o Índice de Facilidade e o Índice de Discriminação ( Ponto-Bisserial ) para cada uma das Questões Objetivas do Componente de Formação Geral. classificadas com Difícil ( 0,16 a 0,40 ) e roxo para as questões classificadas com Muito difícil ( \&lt ; =0,15 ).
Já quanto ao Índice de Discriminação, são usadas as seguintes cores para qualificar a questão : • As questões classificadas com índice Fraco receberam a cor vermelha ( \&lt ; =0,19 ), as classificadas com Médio receberam a cor amarela ( 0,20 a 0,29 ), as classificadas com Bom receberam a cor verde ( 0,30 a 0,39 ) e as classificadas com Muito bom ( \&gt ; =0,40 ) receberam a cor azul.
As Questões Objetivas do Componente de Formação Geral, segundo o Índice de Facilidade, são assim avaliadas : nenhuma das oito questões foi anulada pela Comissão Assessora de Área.
Nenhuma tem o Índice de Facilidade classificado como Muito fácil, e duas questões são tidas como Fácil, por terem índice de acertos iguais a 0,62.
Três questões são consideradas com índice de dificuldade Médio, situando -se no intervalo entre 0,41 e 0,60 do Índice de Facilidade, ou seja, há entre 41,0 \% e 60,0 \% de acertos, enquanto três questões são classificadas na categoria Difícil, apresentando 16 \% a 40 \% de acertos.
Nenhuma questão apresentou menos de 15 \% de acertos, razão pela qual não há questão classificada como Muito difícil.
O Índice de Facilidade varia de 0,19 a 0,62, e o de Discriminação, de 0,23 a 0,49.
Cinco questões, quanto ao Índice de Discriminação, são tidas como Muito bom, duas questões são tidas como Bom e uma questão é classificada como Médio.
Na Tabela 6.9a é apresentada a distribuição das questões, levando -se em conta, simultaneamente, a classificação dos dois índices.
As cinco questões que têm Índice de Discriminação Muito bom figuram entre dois níveis de dificuldade Fácil e Médio : duas classificadas na categoria Fácil ( as Questões 6 e 7 ) do Índice de Facilidade e três, na categoria Médio ( Questões 1, 3 e 5 ).
As Questões 1 e 3 apresentaram o maior poder discriminatório, com índice 0,49, e ambas consideradas Médio em termos de Facilidade, com uma proporção de 0,48 e 0,51 acertos, na ordem.
O máximo de acertos é alcançado pelas Questões 6 e 7, com um Índice de Facilidade de 0,62.
A Questão 8 é considerada Difícil, com Índice de Facilidade 0,19.
Seu Índice de Discriminação é o mais baixo dentre as oito questões válidas : Médio ( 0,23 ).
No Gráfico 6.4a, para exemplificar, analisa -se o comportamento da Questão de número 3 de Formação Geral dos estudantes dos cursos de Licenciatura.
Trata -se de uma questão considerada com nível Médio de dificuldade, em relação à Facilidade, e a que obteve, junto com a Questão de número 1 ( ver Anexo IA ), o maior Índice de Discriminação dessa parte da prova.
Neste gráfico, cada uma das cinco curvas representa o percentual de respostas em determinada alternativa da questão, em função do número de acertos dos estudantes nessa parte da prova ( Formação Geral/Múltipla Escolha ), antes de possíveis eliminações pelo critério Ponto-Bisserial.
A curva em verde corresponde à alternativa B, a alternativa correta para essa questão.
Observa -se que, entre os estudantes com menor número de acertos, nessa parte do exame, a situação mais frequente foi a escolha da alternativa E ( em vermelho ), incorreta.
Por exemplo, entre os estudantes que acertaram duas questões, 45,5 \% escolheram a alternativa E ( em vermelho ), 10,4 \% escolheram a alternativa D ( em laranja ), 5,5 \% escolheram a alternativa C ( em preto ), 27,8 \% escolheram a alternativa B ( em verde, gabarito ), 10,5 \% escolheram a alternativa A ( em azul ).
Dentre os que acertaram duas respostas entre as questões de múltipla escolha de Formação Geral, 0,1 \% deixou a questão em branco e 0,1 \% marcou mais de uma alternativa, invalidando a questão.
À medida que o número de acertos aumenta, indicando desempenho melhor nessa parte da prova, aumenta, concomitantemente, a proporção de estudantes que selecionaram a alternativa correta B, atingindo 100 \% para os estudantes com oito acertos.
Essa análise permite verificar como a questão discriminou os grupos de desempenho, justificando -se o alto Índice de Facilidade ( 0,51 ) obtido na questão.
Cumpre notar que não é possível inferir deste gráfico nem o Índice de Facilidade, que seria uma Média da proporção ponderada pela quantidade de estudantes com cada uma das notas, nem o Índice de Discriminação Ponto-Bisserial, por razão equivalente.
No caso extremo, no qual a grande concentração dos acertos dos estudantes fosse abaixo de quatro, o Índice de Facilidade seria obrigatoriamente abaixo de 60,0 \% ( neste exemplo ).
Caso a concentração fosse em seis acertos ou mais, o índice seria obrigatoriamente acima de 90,0 \%.
Os gráficos relativos às demais questões de Formação Geral constam do Anexo IA. 6.2.2 COMPONENTE DE FORMAÇÃO GERAL – HISTÓRIA ( BACHARELADO ) Na Tabela 6.7b, são apresentadas as Estatísticas Básicas relativas às oito questões objetivas válidas do componente da prova que abrange a Formação Geral dos estudantes de Bacharelado.
A Média do Brasil é 51,4.
A menor Média é encontrada na região Norte ( 37,1 ), e a maior, na região Centro-Oeste ( 56,9 ).
As demais Médias são : 45,6 na região Nordeste, 52,6 na região Sudeste e 50,3 na região Sul.
O Desvio padrão do Brasil é 18,9, sendo o maior Desvio padrão encontrado na região Centro-Oeste ( 25,3 ), e o menor, na região Nordeste ( 17,4 ).
Os demais desvios são : 21,8 na região Norte, 18,3 na região Sudeste e 18,8 na região Sul.
As Medianas do Brasil, como um todo, e das regiões Sudeste e Sul são iguais a 50,0 ; as regiões Norte e Nordeste têm Mediana abaixo do valor Brasil, 37,5 e 43,8, respectivamente, e a Mediana da região Centro-Oeste é a maior ( 56,3 ).
A nota Máxima 100,0 é alcançada nas regiões Sudeste, Sul e Centro-Oeste, na região Norte alcança 75,0 e na região Nordeste alcança 87,5.
As notas Mínimas ( zero ) são iguais para quase todas as regiões, menos para a região Nordeste ( 12,5 ).
Na Tabela 6.8b, são apresentados o Índice de Facilidade e o Índice de Discriminação ( Ponto-Bisserial ) para cada uma das Questões Objetivas do Componente de Formação Geral.
A convenção das cores é a mesma da Tabela 6.8a.
As Questões Objetivas do Componente de Formação Geral, segundo o Índice de Facilidade, são assim avaliadas : das oito questões válidas, nenhuma teve o Índice de Facilidade classificado como Muito fácil, e quatro questões são tidas como Fácil, por terem índice de acertos entre 0,61 a 0,85.
Uma questão é considerada com índice de dificuldade Médio, situando -se no intervalo entre 0,41 e 0,60 do Índice de Facilidade.
Três questões são classificadas na categoria Difícil, situando -se no intervalo entre 0,16 e 0,40.
Nenhuma questão apresenta menos de 15 \% de acertos, razão pela qual não há questão classificada como Muito difícil.
O Índice de Facilidade varia de 0,21 a 0,74 e o de Discriminação, de 0,31 a 0,45.
Seis das oito questões válidas, quanto ao Índice de Discriminação, são tidas como Muito bom e duas como Bom.
Na Tabela 6.9b, é apresentada a distribuição das questões, levando -se em conta, simultaneamente, a classificação dos dois índices.
As seis questões com Índice de Discriminação Muito bom figuram em três níveis de dificuldade Difícil, Médio e Fácil : quatro classificadas na categoria Fácil ( Questões 1, 3, 6 e 7 ) do Índice de Facilidade, uma, na categoria Médio ( a Questão 5 ) e outra, na categoria Difícil ( a Questão 4 ).
Em particular, as Questões 1 e 3 apresentaram o maior poder discriminatório, com índice 0,45, e são consideradas Fácil em termos de Facilidade, com uma proporção de 0,63 acertos, para ambas.
O máximo de acertos é alcançado pela Questão 7, com um Índice de Facilidade de 0,74.
A Questão 8 é considerada Difícil, com Índice de Facilidade 0,21.
O Índice de Discriminação é Bom para a Questão 8 ( 0,31 ).
No Gráfico 6.4b, para exemplificar, analisa -se o comportamento da Questão de número 3 de Formação Geral.
Trata -se de uma questão considerada com nível Fácil de dificuldade, e a que obteve, junto com a Questão de número 1 ( ver Anexo IB ), o maior Índice de Discriminação ( 0,45 ) dessa parte da prova.
Neste gráfico, cada uma das cinco curvas representa o percentual de respostas em determinada alternativa da questão, em função do número de acertos dos estudantes nessa parte da prova ( Formação Geral/Múltipla Escolha ), antes de possíveis eliminações pelo critério Ponto-Bisserial.
Como não houve questões anuladas pela Comissão Assessora de Área o número de acertos alcança o máximo, oito, para esse componente.
A curva em verde corresponde à alternativa B, a alternativa correta para essa questão.
Observa -se que entre os estudantes com menor número de acertos, nessa parte do exame, a situação mais frequente foi a escolha da alternativa E ( em vermelho ), incorreta.
Por exemplo, entre os estudantes que acertaram duas respostas, 44,7 \% escolheram a alternativa E ( em vermelho ), 8,0 \% escolheram a alternativa D ( em laranja ), 8,7 \% escolheram a alternativa C ( em preto ), 24,0 \% escolheram a alternativa B ( em verde, gabarito ), 14,7 \% escolheram a alternativa A ( em azul ).
Dentre os que acertaram duas respostas entre as questões de múltipla escolha de Formação Geral, nenhum deixou a questão em branco ou marcou mais de uma alternativa, invalidando a questão.
À medida que o número de acertos aumenta, indicando desempenho melhor nessa parte da prova, aumenta, concomitantemente, a proporção de estudantes que selecionaram a alternativa correta B, atingindo 100 \% para os estudantes com oito acertos.
Essa análise permite verificar como a questão discrimina os grupos de desempenho, justificando -se o alto índice ( 0,45 ) obtido na questão.
Cumpre notar que não é possível inferir deste gráfico nem o Índice de Facilidade, que seria uma Média da proporção ponderada pela quantidade de estudantes com cada uma das notas, nem o Índice de Discriminação Ponto-Bisserial, por razão equivalente.
No caso extremo, no qual a grande concentração dos acertos dos estudantes fosse abaixo de quatro, o Índice de Facilidade seria obrigatoriamente abaixo de 60,0 \% ( neste exemplo ).
Caso a concentração fosse em seis acertos ou mais, o índice seria obrigatoriamente acima de 80,0 \%.
Os gráficos relativos às demais questões de Formação Geral constam do Anexo IA. 6.2.3 COMPONENTE DE CONHECIMENTO ESPECÍFICO – HISTÓRIA ( LICENCIATURA ) Na Tabela 6.10a, são apresentadas as Estatísticas Básicas em relação às Questões Objetivas do Componente de Conhecimento Específico da prova de História ( Licenciatura ) por Grande Região.
Nessa parte da prova, como será discutido ainda nesta seção, cinco questões não serão usadas no cômputo das notas por terem sido descartadas pelo critério do Ponto-Bisserial.
Assim, as notas foram calculadas com base em 22 das 27 Questões Objetivas válidas de Conhecimentos Específicos.
A Média do Brasil deste componente é 47,8.
A menor Média é observada na região Sul ( 44,7 ), e a maior, na região Sudeste ( 50,5 ).
O Desvio padrão de todo o Brasil é 19,9, sendo o menor Desvio padrão encontrado na região Norte ( 18,5 ), e o maior, na região Sudeste ( 20,5 ).
A Mediana de todo o Brasil é 45,5, o mesmo valor da Mediana encontrado nas regiões Norte, Nordeste e Centro-Oeste.
Na região Sudeste ( 50,0 ) o valor é maior, e, na região Sul ( 40,9 ), o valor é menor.
A nota Máxima da prova é 100,00, obtida nas Questões Objetivas do Componente de Conhecimento Específico, por, pelo menos, um estudante das regiões Nordeste, Sudeste e Sul.
Na região Centro-Oeste o valor é 95,5 e a menor nota Máxima ocorre na região Norte ( 90,9 ).
A nota Mínima é zero, obtida, por pelo menos, um estudante em todas as regiões nas Questões Objetivas do Componente de Conhecimento Específico.
Na Tabela 6.11a, são apresentados os Índices de Facilidade e de Discriminação ( Ponto-Bisserial ) das Questões Objetivas do Componente de Conhecimento Específico, para os estudantes de Licenciatura em História.
Para facilitar a diferenciação das questões, serão usadas as mesmas cores da Tabela 6.8a e 6.8b para as diferentes classificações dos Índices de Facilidade e de Discriminação.
Sendo assim, as análises a seguir dizem respeito às 27 questões válidas.
A partir do Índice de Facilidade obtido, conclui -se que quase metade das questões objetivas da prova são consideradas Difícil : das 27 questões válidas, 13 são classificadas como Difícil.
Onze questões foram classificadas como Médio, e outras três consideradas como Fácil.
Nenhuma questão foi classificada como Muito fácil ou Muito difícil.
Já quanto aos Índices de Discriminação das Questões Objetivas do Componente de Conhecimento Específico da prova, obtêm -se como resultado a seguinte classificação : sete das 27 questões válidas são consideradas boas, e outras 13 têm Índice de Discriminação Muito bom.
Assim, para 20 em 27 questões, os Índices de Discriminação são Bom ou Muito bom.
Dentre as demais, duas delas são classificadas como Médio, e outras cinco, como Fraco, sendo sete, por conseguinte, a quantidade de questões nos dois patamares mais baixos de Discriminação.
Constata -se, assim, que a prova – no que se refere ao Componente de Conhecimento Específico – possuía boa capacidade de discriminar entre aqueles que dominam ou não o conteúdo.
O Índice de Facilidade variou de 0,27 a 0,69 e o de Discriminação, de 0,03 a 0,55.
Na Tabela 6.12a, é apresentada a distribuição das questões levando -se em conta, simultaneamente, a classificação dos dois índices.
Dentre as questões que alcançam os maiores Índices de Discriminação, a Questão 15 é classificada com Índice de Discriminação Muito bom.
Seu índice é 0,55, e, quanto ao Índice de Facilidade, é classificada como Médio, com 48,0 \% dos estudantes marcando a opção correta.
A Questão de número 26 foi a mais difícil entre as 27 questões específicas válidas, com baixo Índice de Facilidade, apenas 27,0 \% de acertos.
Essa questão apresenta poder discriminatório muito baixo, 0,08, o que comprova ter sido a mais difícil para os estudantes.
Destaca -se, também, a Questão 24, com Índice de Facilidade 0,37, o que, em termos percentuais, corresponde a 37,0 \% de estudantes que responderam acertadamente, sendo, também, 0,03 seu Índice de Discriminação.
Tais questões foram, portanto, pelo critério Ponto-Bisserial, consideradas inadequadas.
Por isso, as Questões 24 e 26, além das Questões 9, 11 e 30, foram eliminadas do cômputo da nota final.
A título de exemplo das análises do comportamento das questões objetivas, no Gráfico 6.5a, analisa -se a Questão 15 do Componente de Conhecimento Específico.
Essa questão é considerada, pelas respostas dos estudantes avaliados na prova, como uma questão com nível de Facilidade Médio, 0,48, ou seja, 48,0 \% dos estudantes assinalaram, acertadamente, a opção C ( linha preta ) correspondente ao gabarito.
Como já comentado, seu Índice de Discriminação é igual a 0,55, classificado como Muito bom, o maior valor de Discriminação.
Neste gráfico, cada uma das cinco curvas representa o percentual de respostas em determinada alternativa da Questão 15, em função do número de acertos dos estudantes nessa parte da prova, antes de possíveis eliminações de questões pelo critério Ponto-Bisserial.
A alternativa correta C, representada no gráfico pela curva em preto, foi escolhida em maiores proporções pelos estudantes com desempenho melhor nessa parte da prova.
Já as alternativas incorretas, também denominadas distratores, foram selecionadas, principalmente, por aqueles com notas mais baixas.
No caso de História ( Licenciatura ), como se observa no eixo horizontal do Gráfico 6.5a, nenhum estudante acerta mais do que 25 questões dentre as 27 questões objetivas válidas, considerando -se o gabarito original.
A eliminação das cinco questões com Índice de Discriminação Fraco do cômputo da nota é que possibilitou a existência de nota Máxima igual a 100,0 por pelo menos um estudante das regiões Nordeste, Sudeste e Sul ( visto na Tabela 6.10a ).
Observa -se que a soma não é 100 \%, por causa das questões não respondidas ou com mais de uma opção marcada.
Aqueles com nota zero, na sua quase totalidade, marcaram a opção B ( linha verde ).
A proporção de estudantes que selecionaram a resposta correta C aumenta gradativamente, chegando a atingir 100 \% para 23 acertos ou mais, enquanto a proporção dos que escolheram alternativas incorretas decai, a partir de sete acertos, em função do número de acertos nessa parte da prova.

% matched lemmas: 
\end{textsample}
